
%-------------------------------------------------------------------
% Краткое руководство по базе графического и цифрового дизайна.
% Информация актуальна на август 2026. Используетcя модифицированный
% пакет modernclassnotes~"--- скачайте в репозитории. Это временное
% решение~"--- по ходу написания переедем на другой пакет, либо
% заново переопределим цветовую гамму акцентов. 
%-------------------------------------------------------------------

\documentclass[%
	mode=chapternotes,
	palette=mono,
	fontsize=14pt,
	official=false,
]{preamble/modernclassnotes}

% Перенесено в classnotes.sty
% \usepackage[T2A]{fontenc}

% Делаем с main, чтобы подгрузились настройки переносов и прочего
\usepackage[english, main=russian]{babel}

\usepackage{pscyr}

% -------------------------------------------------------------------
% Переопределение для корректной работы пакета modernclassnotes
% -------------------------------------------------------------------

\renewcommand{\rmdefault}{ftx}
\renewcommand{\sfdefault}{ftx}
\renewcommand{\familydefault}{\sfdefault}

% ======================== Микротипография ==========================

\usepackage[%
	kerning=true,
	spacing=true,
	protrusion=true,
	babel=true,
	factor=1000
]{microtype}

% ============================ minted ===============================

\usepackage{minted}

\setminted{
	style=friendly,
	linenos,
	numbersep=8pt,
	breaklines,
	autogobble,
	tabsize=4,
	bgcolor=mcnBg
}

\usepackage{graphicx}

% https://tex.stackexchange.com/a/617816
% \makeatletter
% \graphicspath{{\input@path/img/}}
% \makeatother


% Новые окружения
% \begin{...} ... \end{...}
\newminted{html}{}  % htmlcode
\newminted{latex}{} % latexcode

% Новые команды для внутристрочных вставок
\newmintinline[htmlinline]{html}{}   % \htmlinline{...}
\newmintinline[latexinline]{latex}{} % \latexinline{...}

% vim:tw=69


\begin{document}

\course{Цифровой и продуктовый дизайн}

\docsubtitle{Введение и~источники для~самостоятельной учёбы}

\title{Цифровой и~продуктовый дизайн}

\maketitle
\tableofcontents
\newpage

% -------------------------------------------------------------------
% 1. Инструменты и вводная
% -------------------------------------------------------------------

\lecture[17 августа 2026]{1}{Инструменты и теория}

\textbf{Дисклеймер.} В~этой инструкции описаны основные определения
и~собраны ссылки на~материалы для~обучения в~помощь начинающим
дизайнерам. Все~книги и~курсы авторы читали и~проходили
самостоятельно, поэтому информация предвзята, субъективна
и~не~является актуальным отражением положения дел в~индустрии.
На~момент написания авторы находились на~джуниор"=миддл"=позиции,
поэтому, пожалуйста, относитесь ко~всей информации ниже скептически
и~ищите дополнительные источники информации. Поиск и~проверка
информации~"--- ваш базовый навык как профессионала в~любой сфере,
особенно актуальный во~времена галлюционирующего ИИ~:-)

\begin{definition}[кто такой продуктовый дизайнер]
	\textbf{Продуктовый дизайнер}~"--- человек, использующий
	инструменты и~методологию дизайна при~работе над~цифровым
	продуктом и~влияющий на~его~развитие. 
\end{definition}

Продуктовый дизайнер отвечает за то, что~происходит в~продукте:
проектирует новые фичи, реагирует на~обратную связь пользователей
и~бизнеса [Дополнить].  Ключевый термин в~работе продуктового
дизайнера~"--- \textbf{метрики}, или~числовые показатели продукта.
Существует большое количество метрик для~оценки состояния
продукта~"--- например, MAU, DAU, Retention, CTR [Вставить ссылку],
о~которых мы~поговорим дальше.

\begin{warning}
	Не все школы дизайна считают, что~продуктовый дизайнер обязан
	погружаться в~метрики. Тем не менее, полезно знать хотя бы
	некоторые из~них, чтобы говорить на~одном языке с~менеджерами
	продукта и~стейкхолдерами~"--- владельцами бизнеса. 
\end{warning}

Чтобы стать дизайнером, рекомендуем начать по~порядку.

% Делаем и элементы списка полужирным. — Rikudo
\begin{enumerate}[{\bfseries 1.}]
		\item \textbf{Инструменты:} Фигма, Клод, Клод Код, ЧатЖПТ,
			Миджорни (базово). На~август 2026 rода рекомендуем добавить
			Гитхаб и~по~желанию~"--- дополнительные инструменты работы
			с~графикой (Фотошоп, Иллюстратор, Блендер, Афтер~Эффектс).

		\item \textbf{Теория:} типографика, цвет, композиция,
			интерфейсы.
\end{enumerate}

%\begin{takeaways}
%\begin{itemize}
%   \item Вывод «А»;
%   \item Вывод «Б»;
%   \item Вывод «В»;
%\end{itemize}
%\end{takeaways}

% Лучше \clearpage. — Rikudo
% \newpage
\clearpage


% -------------------------------------------------------------------
% 2. Типографика~"--- определения и источники
% -------------------------------------------------------------------
\lecture[2026]{2}{Типографика}


% В кодбокс не вставляется русский текст~"--- поправить
%
% Это проблема пакета lstlisting. Он плохо работает с кириллицей, так что
% предлагаю перейти на пакет minted
\begin{codebox}[%
		Базовый HTML~"--- чтобы разбираться, как устроены веб-страницы
		с~помощью инспектора элементов%
	]
	\inputminted{html}{code/html_1.html}
\end{codebox}


%\begin{tip}
%Совет, полезная информация
%\end{tip}

%\begin{solution}
%Решение чего-либо
%\end{solution}


\end{document}

% vim:ts=2:sw=2:tw=68
